%% kapa: exam=1 round=22 year=2011 period=2 form=A subject=appraisal_law subject_ko=부동산관계법규 seq=2 number=42
%% source: Q-Net 감정평가사 기출문제 (cst00309) artlSeq=5036790 file=extracted/1008/5036790/2011년 제22회 감정평가사 1차시험 기출문제/2교시.hwp
%% conversion: kordoc 4.24.51 -> tools/gen1.py (auto); review=proofread-gpt-astra; flags=none
\begin{question}{42}
국토의 계획 및 이용에 관한 법령상 광역도시계획에 관한 설명으로 옳은 것은?
\begin{choices}
  \item 중앙행정기관의 장, 특별시장ㆍ광역시장ㆍ도지사ㆍ특별자치도지사, 시장 또는 군수는 국토해양부장관이나 도지사에게 광역계획권의 변경을 요청할 수 있다.
  \item 광역계획권이 둘 이상의 시ㆍ도의 관할 구역에 걸쳐 있는 경우에는 관계 시ㆍ도지사가 협의하여 광역계획권을 지정한다.
  \item 광역계획권이 같은 도의 관할 구역에 속하여 있는 경우에는 도지사가 광역도시계획을 수립하여야 한다.
  \item 도지사가 광역계획권을 지정하거나 변경하려면 국토해양부장관과의 협의와 관계 시ㆍ도, 시 또는 군의회의 의견을 들은 후 지방도시계획위원회의 심의를 거쳐야 한다.
  \item 광역도시계획을 공동으로 수립하는 시ㆍ도지사는 그 내용에 관하여 서로 협의가 되지 아니하면 공동으로 국토해양부장관에게 조정을 신청하여야 한다.
\end{choices}
\end{question}
\answer{1}
